%% kapa: exam=1 round=26 year=2015 period=1 form=A subject=civil_law subject_ko=민법 seq=40 number=40
%% source: Q-Net 감정평가사 기출문제 (cst00309) artlSeq=5138001 file=extracted/1008/5138001/제26회_감정평가사 1차_기출/제26회_감정평가사 1차_기출/1교시_감평_A형.pdf
%% conversion: kordoc 4.24.51 -> tools/gen1.py (auto); review=proofread-sonnet; flags=none
\begin{question}{40}
저당권에 관한 설명으로 옳지 않은 것은? (다툼이 있으면 판례에 따름)
\begin{choices}
  \item 저당권설정자는 현재 저당부동산의 소유자가 아니라면 피담보채무가 소멸하더라도 저당권의 말소를 청구할 수 없다.
  \item 저당권은 경매에서의 매각으로 인하여 소멸한다.
  \item 저당권은 피담보채권과 분리하여 타인에게 양도하거나 다른 채권의 담보로 하지 못한다.
  \item 저당권자가 물상대위를 통하여 우선변제를 받기 위해서는 저당권설정자가 받을 가치 변형물을 그 지급 또는 인도 전에 압류하여야 한다.
  \item 저당권자가 물상대위권을 행사하지 아니한 경우, 저당목적물의 변형물로부터 이득을 얻은 다른 채권자에 대하여 부당이득반환을 청구할 수 없다.
\end{choices}
\end{question}
\answer{1}
